% Introduction: API loop, addressing convention, and the variable taxonomy.

The MODFLOW 6 Application Programming Interface (API) provides access to the
program's internal memory while a simulation is running. The API implements the
Basic Model Interface (BMI) and an eXtended Model Interface (XMI) that exposes
the nonlinear solution loop. The API is distributed
as a shared library (\texttt{libmf6}) and is most commonly driven from Python
through the \texttt{modflowapi} and \texttt{xmipy} packages.

Every variable stored in the MODFLOW 6 memory manager is exposed through the
API. A variable is referenced by its \emph{memory address}, which has the form

\begin{quote}
\texttt{MODELNAME/PACKAGENAME/VARNAME}
\end{quote}

\noindent where all components are uppercase (for example,
\texttt{MYMODEL/NPF/CONDSAT} or \texttt{MYMODEL/STO/SS}).

\subsection{Reading and Writing Variables: Pointers versus Copies}

A variable's value is accessed through one of three routines, which differ in
both cost and semantics. Choosing the right one matters for performance in a
time loop and for correctness when a value is written.

\begin{description}
\item[\texttt{get\_value\_ptr}] returns an array that is a direct view of
MODFLOW 6's memory; no data are copied. Reading it always reflects the current
value, and assigning into it in place writes directly into MODFLOW 6's memory.
Because nothing is copied, a pointer can be obtained once (for example, just
after \texttt{initialize}) and reused for the entire simulation, which makes it
the most efficient option for a variable accessed every time step or every
iteration. The view shares MODFLOW 6's exact data type and shape, so it must be
written with the matching type and dimensions. A pointer is valid only while the
underlying array remains allocated in place; if MODFLOW 6 reallocates the
variable the pointer becomes stale and must be obtained again. Writing through a
pointer bypasses the memory-set handler described below, and pointers are not
available for string variables.

\item[\texttt{get\_value}] returns a \emph{copy} of the value. The copy is
independent of MODFLOW 6's memory: it is a snapshot that does not change when
MODFLOW 6 later updates the variable, and it cannot be invalidated by
reallocation. Each call allocates and fills a new array, costing memory and time
proportional to the variable's size, so it is best for occasional reads or when
an unchanging snapshot is wanted (for example, to compare a value before and
after a solve).

\item[\texttt{set\_value}] copies a supplied array \emph{into} MODFLOW 6's
memory. Like \texttt{get\_value} it copies element by element and pays that cost
on every call. Unlike a write through a pointer, \texttt{set\_value} invokes
MODFLOW 6's memory-set handler, so it is the appropriate choice if a variable
ever triggers a recomputation when it is set externally. No standard package
registers such a handler at present, but the mechanism exists and
\texttt{set\_value} is the forward-compatible setter.
\end{description}

\noindent In short: obtain a pointer with \texttt{get\_value\_ptr} and reuse it
when repeatedly reading or writing the same variable inside the time loop; use
\texttt{get\_value} when an independent snapshot is wanted; and use
\texttt{set\_value} for a single copy-in or when honoring a memory-set handler
matters. The pseudo-code and example in this report obtain a pointer once and
write through it, which is the common pattern for an API time loop. The named
routines follow the \texttt{xmipy} and \texttt{modflowapi} interfaces; the
\texttt{modflowapi} package additionally offers higher-level helpers, but the
cost and semantics described here are the same.

\subsection{The Simulation Loop and When Variables Are Written}

Whether a modification ``takes'' depends entirely on \emph{when} MODFLOW 6 next
writes the variable relative to when the user sets it. A MODFLOW 6 simulation
advances through a fixed sequence of phases, and each package writes its
memory-manager variables at well-defined points in that sequence:

\begin{itemize}
\item \textbf{Allocate and Read (AR)} --- executed once, after input is read.
Many derived quantities (for example, saturated conductance in the Node
Property Flow Package) are computed here from the raw input arrays and then used
unchanged for the rest of the simulation.
\item \textbf{Read and Prepare (RP)} --- executed at the start of each stress
period. List-based and array-based stress data (for example, well pumping rates)
are re-read from the input context here, overwriting any externally set values.
\item \textbf{Advance (AD)} --- executed at the start of each time step.
Time-varying inputs (for example, time-varying hydraulic conductivity, TVK) and
the quantities derived from them are refreshed here when those options are
active.
\item \textbf{Formulate (CF/FC/FN)} --- executed at every outer (Picard or
Newton) iteration. Variables read directly here (for example, specific storage
and specific yield in the Storage Package) reflect whatever value is currently
in memory at each iteration.
\item \textbf{Calculate Flows (CQ)} and output --- executed at the end of each
time step.
\end{itemize}

\noindent When using \texttt{modflowapi}, these phases correspond to callback
hooks (such as the stress-period and time-step callbacks); a value should be set
in the callback that fires \emph{after} MODFLOW 6 last writes the variable and
\emph{before} the variable is next used. As a general rule, values are set after
\texttt{prepare\_time\_step} and before the solve.

\subsection{Variable Modifiability Classes}

Each variable documented in this report is assigned one of the following
classes. The class can depend on which package options are active, as noted in
the per-variable entries.

\begin{description}
\item[Live input] An input array that is read once (or re-read each stress
period) and then used directly in the formulate step. Setting it through the API
changes the simulation, and the change persists until MODFLOW 6 next reads that
input. Example: \texttt{STO/SS}, \texttt{STO/SY}.

\item[Derived variable] A quantity computed from input during AR and then used
directly thereafter. This is the variable to set in order to change the
associated behavior; setting the raw input that produced it has no effect.
Example: \texttt{NPF/CONDSAT} (set this to change intercell conductance).

\item[Inert input] An input that is used only to compute a derived variable
during AR. Setting it after AR has no effect, because the derived quantity has
already been computed. To achieve the intended change, set the corresponding
derived variable instead. Example: \texttt{NPF/K} (when time-varying
conductivity is not active).

\item[Period-reset] A variable that is re-read from the input context at the
start of every stress period. It can be set through the API, but the value is
overwritten at the next Read and Prepare. Example: well pumping rates in
\texttt{WEL/BOUND}.

\item[Overwritten] A variable that MODFLOW 6 recomputes every time step (or
every iteration). Setting it has no lasting effect. This class includes derived
knobs whose source inputs are refreshed each time step when an option such as
TVK, time-varying storage (TVS), the Viscosity Package (VSC), or a time series
is active.

\item[State] Solution state owned by the solver or model, such as the dependent
variable (head). These can be set through the API but require care, because they
participate directly in the nonlinear solution.

\item[Structural / read-only] Grid geometry, cell connectivity, and array sizes.
These define invariants relied upon throughout the program and must be treated
as read-only. Example: \texttt{DIS/NODES}, connection arrays.
\end{description}

\subsection{How to Read the Package Tables}

Each package section lists the memory-manager variables that an API user is
likely to want to modify. For each variable the table gives its name (the
\texttt{VARNAME} component of the memory address), its modifiability class, and
guidance describing what to set, when to set it, any options that change the
classification, and the variable to set instead when the listed variable is
inert or overwritten. Variables that are purely structural or that are internal
bookkeeping are generally omitted unless they are a common source of confusion.
